\section{System Description}
\label{sec:model}

%In this section, we describe our system model.
We consider a single server queueing system with two job (a.k.a.
customer) classes: eager customers (also denoted as
$\epsilon$-customers) demand service immediately upon arrival, whereas
tolerant customers (also denoted as $\tau$-customers) can wait in a
queue (of infinite capacity) to be served.\footnote{\revision{We
    discuss the generalization where eager customers have limited
    patience and abandon/renege after a short wait time in
    Section~\ref{sec:conclusion}; see also~\cite{ANOR}.}} The
$\tau$-customers can be interrupted either partially (i.e., their
service rate may be reduced) or completely by $\i$-customers, but not
by other $\tau$-customers.  Without loss of generality, we assume {\it
  a unit server speed.}
%
We assume that $\epsilon$-customers (respectively, $\t$-customers)
arrive according to a Poisson process with rate $\lambda_{\i}$
(respectively, $\lambda_\t$). The sequence of job sizes
(a.k.a. service requirements) for both the classes is i.i.d., with
$B_\i$ denoting a generic $\i$ job size, and $B_\t$ denoting a generic
$\t$ job size. Throughout, {\it we assume that $B_\t$ is exponentially
  distributed with mean $1/\mu_\t,$ and that $\Exp{B_\i^2} < \infty.$}
Let $\mu_\i := 1/\Exp{B_\i}.$

\subsection{Dynamic schedulers}

We consider dynamic scheduling, wherein the scheduling policy for each
class depends upon the state (occupancy) of both the classes.  An
eager scheduling (sub)policy (a set of rules that assign fractions of
server capacity to various eager customers, which may also include the
admission control rules) is chosen depending upon the tolerant
state. The overall dynamic scheduler is characterized by a sequence of
such eager sub-policies, one corresponding to each value of tolerant
occupancy---the tolerant queue in turn utilizes the service capacity
left unused by the eager class in a work-conserving manner. Our
dynamic schedulers can thus be viewed to be of nested type: a
top-level policy chooses the sub-policy used for scheduling the
$\i$-class based on the occupancy (state) of the
$\t$-class. \revision{The sub-policy in turn determines, based on the
  number of eager customers in the system, the admission rule for
  subsequent arrivals, and the service rate to allocate to each
  existing eager job. It is important to note that these nested
  schedulers are not restrictive; rather this is simply a convenient
  representation of dynamic schedulers; details and examples follow.}
As a result, the service processes of the two classes are
interdependent (unlike in the case of static scheduling as considered
in \cite{ANOR,Value}).

\revision{Specifically, let $(X_\t (t), X_\i (t))$ represent the
  number of tolerant and eager customers in the system at time~$t$. We
  consider stationary Markov scheduling policies that determine the
  service plan for both classes (including, possibly, admission
  control of eager customers), depending on the tuple
  $(X_\t (t), X_\i (t)).$
  % With new arrival/departure of $\t$ or $\i$ class, the
  % scheduling policy changes the service plan of both the classes.
  Such Markov policies are known to be sufficient for a wide range of
  sequential decision problems (see \cite{Puter}); in the context of
  our model, they are sufficient if one additionally assumes that
  eager customers have exponentially distributed service requirements.
  % We consider stationary Markov
  % scheduling
  % policies, as it is well known that these policies are sufficient
  % for
  % a wide range of sequential decision problems (\cite{Puter}); when
  % one considers exponential service times for eager class also, then
  % these policies are sufficient for the problems considered in
  % section
  % \ref{sec:pareto}.
  % For example, if $X_\t (t) = j$ the $\i$-class is served according
  % to a given sub-policy $j$ and the remaining server capacity (which
  % depends upon the number $\i$-customers) is used to serve
  % $\t$-class.
  As an example, consider the following policy parameterized by
  $(L, N) \in \mathbb{N}^2$: a) when $(X_\t(t), X_\i (t)) = (j, c)$
  with $c \le N$, and $j \le L$ then each $\i$-customer is served with
  rate (speed) $1/N$, while longest waiting $\t$-customer is served at
  rate $(N-c)/N$; and b) if $j > L$, the longest waiting $\t$-customer
  is served with the full service capacity (i.e., at rate 1), and no
  further $\i$-customers are admitted. It is convenient to view such
  policies as being \emph{nested}; for the above example, when
  $X_\t (t) \le L$ the eager class is served according to one
  sub-policy, while it is served according to a second sub-policy when
  tolerant occupancy is more than $L$. Under the former sub-policy,
  the eager class is served as an $M/G/N/N$ queue (with each `server'
  taken to have speed~$1/N$), while under the latter sub-policy, no
  eager customers are admitted. There is of course the issue of how to
  deal with the existing eager customers in the system when there is
  an arrival/departure in the $\t$ queue, leading to change in the
  eager sub-policy. In the above example, this corresponds to the
  handling of any eager customers that remain when the $\t$ occupancy
  increases from $L$ to $L+1.$ This issue is addressed later in this
  section (see Assumption {\bf A}.1 and the discussion in
  Section~\ref{sec:example_models}).}

Note that the eager sub-policy switches whenever there is an
arrival/departure in the tolerant queue, though the further details of
each sub-policy depend only on the eager occupancy.
%the eager scheduling policy also depends upon eager occupancy;
%the scheduling policy of the eager class is (possibly) changed only
%when the tolerant state changes, i.e., eager class can be scheduled
%using new rules, when the tolerant state changes.
We typically refer to the $\i$-sub-policy implemented when there are
$j$ tolerant jobs in the system as sub-policy~$j.$

\noindent {\bf $\i$-schedulers:} Note that while the occupancy of the
tolerant queue dictates the selection of the eager sub-policy, the
sub-policies themselves are oblivious to the state of the tolerant
queue.
%Moreover, what we refer to as a sub-policy includes system
%decisions (e.g., service capacity allocated to the $\i$-class,
%admission control, amount of waiting space, etc.) as well as
%behavioural aspects of the impatient eager customers (e.g., eager
%customers may balk based on the system occupancy).
%
We make the following additional assumptions.  Some examples of
schedulers that satisfy these are presented in
Section~\ref{sec:example_models}:
\begin{enumerate}[{\bf A}.1]
\item To simplify the transition from one $\i$-sub-policy to the next,
  we assume that all $\i$-customers are dropped when there is an
  arrival/departure in the $\t$-queue.
%  \footnote{Assumption {\bf A}.1
%    is a technical condition required in our proofs. However, under
%    the SFJ limit, this `flushing' of the $\i$-system is only
%    performed at a bounded rate (since arrivals/departures in the
%    $\t$-queue occur at a bounded rate), while the arrival rate of the
%    $\i$ system scales to infinity. Thus, we expect that this
%    assumption will not impact the blocking probability of the eager
%    class (this is also evident from the Monte Carlo simulation based
%    study presented in Section~\ref{sec:pareto}).}
\item The scheduling of each sub-policy depends only on the number of
  $\i$-jobs present in the system.
\item On arrival, an eager job is either admitted into service or gets
  dropped/blocked. If admitted, an eager job receives service at a
  minimum rate/speed of $c_{\min} > 0$ until its departure \emph{under
    any sub-policy.}
  % \sout{There exists a uniform upper bound $\mathcal{K}$ on the
  % number of $\i$-jobs in the system at any time across
  % sub-policies.}  
%\item Across sub-policies, the second moment of the $\i$-busy cycle,
%  defined as the interval between the start of two successive $\i$
%  busy periods, is uniformly bounded.
\end{enumerate}
%

Assumption {\bf A}.1 is a technical condition required in our proofs.
In Section~\ref{sec:pareto}, we show that this assumption has
negligible influence on the performance of either class under our
partial fluid limit.\footnote{ Under the SFJ limit, this `flushing' of
  the $\i$-system is only performed at a bounded rate (since
  arrivals/departures in the $\t$-queue occur at a bounded rate),
  while the arrival rate of the $\i$ system scales to infinity. Thus,
  we expect that this assumption will not impact the blocking
  probability of the eager class (this is also evident from the Monte
  Carlo simulation based study presented in
  Section~\ref{sec:pareto}).} \revision{We are also able to relax
  Assumption {\bf A}.1 for the special case of exponentially
  distributed eager jobs (see Theorem~\ref{Thm_withou_A1}). As
  discussed before, Assumption {\bf A}.2 is not restrictive; it allows
  all (non-size based) Markov policies.
  % This assumption allows
  % a broad class of schedulers even for the general case,
  We provide examples of such schedulers in
  Section~\ref{sec:example_models}.}
% 
Assumption~{\bf A}.3 implies that the eager class operates as a
\emph{loss system}. Indeed, the requirement of a minimum service
rate/speed $c_{\min} > 0$ captures the `eagerness' of eager
customers. This assumption also implies that there exists a uniform
upper bound $\mathcal{K}$ on the number of $\i$-jobs in the system at
any time \emph{across sub-policies} (since the server is assumed to
have a unit speed).

The above assumptions imply the following condition (proved as
Lemma~\ref{lemma:bp_upper_bound} in Appendix~\ref{sec:busycycles});
the same is stated as a (redundant) assumption to emphasize that the
uniform convergence required in our analysis is predominantly due to
the following uniform bound on the eager busy cycles:
\begin{enumerate}[{\bf A}.4]
\item Across sub-policies, the second moment of the $\i$-busy cycle,
  defined as the interval between the start of two successive $\i$
  busy periods, is uniformly bounded.
\end{enumerate}
%

%Assumption~{\bf A}.4 may appear restrictive at first glance, but is
%actually a consequence of Assumptions~{\bf A}.1--3.  (see
%Lemma~\ref{lemma:bp_upper_bound} in
%Appendix~\ref{sec:busycycles}). Despite its redundancy, we state it
%here to highlight the key role it plays in our analysis.
%


 

% It is important to note that the above assumptions do not impose
% that eager customers must begin service upon arrival. Instead,
% Assumptions~{\bf A}.1--4 allow for a rich class of models for the
% scheduling and behaviour of the eager class, including balking and
% reneging (see Section~\ref{sec:example_models}). Indeed, the
% `eagerness' of the eager class is captured by Assumptions~{\bf A}.3
% and~{\bf A}.4, which impose uniform (across $\i$ sub-policies) upper
% bounds on (i) the number of eager jobs in the system at any time,
% and (ii) the second moment of eager busy cycles. Assumption~{\bf
% A}.4 can usually be shown to hold so long as
% $\Exp{B_\i^2} < \infty$ by upper bounding by the eager busy period
% by the busy period of an associated $M/G/\infty$ queue, which is a
% well understood object \cite{??}; see
% Section~\ref{sec:example_models}.

\noindent {\bf $\t$-schedulers:} Next, we state our assumptions on the
scheduling policy of the tolerant class.
\begin{enumerate}[{\bf B}.1]
\item The $\tau$-scheduler is work conserving, i.e., it utilizes all
  the service capacity left unused by $\epsilon$-jobs, so long as the
  $\tau$-queue is non-empty.
\item The $\tau$-jobs are served in a serial fashion, i.e.,
  $\tau$-jobs cannot pre-empt one another.
\item The $\tau$-scheduler is blind to the size of $\tau$-jobs.
\end{enumerate}

Assumption~{\bf B}.1 implies that the tolerant class experiences a
time varying service process, which depends on both the $\tau$-state
as well as the $\epsilon$-state. Assumptions~\textbf{B}.2-3 imply that
we consider $\tau$-schedulers which are non-pre-emptive and
non-anticipative, for instance, {\it first come first served} (FCFS),
{\it last come first served} (LCFS), and {\it random order of service}
(see
\cite[Chapter~29]{Harchol-Balter}). %Some example models are described
%at the end of this section.

We require another assumption ({\bf B}.4) regarding the stability of
the $\t$-queue under the SFJ limit, when $\i$-customers employ a
single sub-policy (irrespective of the $\t$-state).
% These are referred to as $\t$-static schedulers in \cite{ANOR}.
We provide the required background on these schedulers, define
formally the SFJ scaling, and state the resulting pseudo-conservation
law (see \cite{ANOR} for more details), after which we state
Assumption~{\bf B}.4.


\subsection{$\t$-static schedulers and background }
\label{sec:background}

We now consider the special case of $\t$-static scheduling, where a
single $\i$-sub-policy is used at all times (irrespective of
$\t$-state); this case was analysed in~\cite{ANOR}. Let ${P_B}_j$
represent the blocking probability \kt{(long run fraction of losses)}
of the $\i$-class, if sub-policy-$ j$ is used in a $\t$-static manner
(${P_B}_j$ was referred to as the standalone blocking probability of
sub-policy~$j$ in Section~\ref{sec:intro}); we also refer to these as
$\t$-static blocking probabilities.

%
\noindent {\bf Short-Frequent Jobs (SFJ) Scaling:} Under the SFJ
scaling (as in \cite{ANOR}), we let $\lambda_\i \ra \infty$ and
$\mu_\i \ra \infty,$ such that $\rho_\i := {\lambda_\i}/{\mu_\i}$
remains constant. This corresponds to scaling the arrival as well as
the service rate of the eager class to infinity proportionately, so
that the offered load (the long term rate at which work arrives into
the system) is held constant. We use $\mu_\i$ as the scale parameter
for this partial scaling.
%
Specifically, we scale the job size distribution of the eager class
as, $B^{\mu_\i}_\i \eqdist B^1_\i/\mu_{\i}$, where $B^{\mu_\i}_\i$
denotes a generic eager job size at scale $\mu_\i$ and $\eqdist$ is
equality in distribution. This scaling (plus Poisson arrivals) under
{\bf A}.2 ensures that at scale $\mue,$ the occupancy process of the
$\i$-class gets time-scaled (fast-forwarded) by~$\mue;$ the details
can be found in Appendix~A. Note that the tolerant workload remains
unscaled. Thus, the SFJ scaling may be viewed as a time-scale
separation, with the eager class operating at a faster time-scale.

%
\noindent {\bf Static Pseudo Conservation:} Let
$\Omega^{\mu_\epsilon}_{j}(t)$ represent the total amount of server
capacity left unused by the $\epsilon$-customers in time interval
$[0,t]$, under sub-policy $j$ operating in a $\t$-static manner. Note
that $\Omega^{\mu_\epsilon}_{j}(t)$ is the (cumulative) service
process seen by the $\t$-system. Then by \cite[Lemma 1]{ANOR}, for all
$\mue,$ the asymptotic (in time) growth rate of
$\Omega^{\mu_\epsilon}_{j}(t)$ satisfies:

\vspace{-5mm}
{\small \begin{align}
    \label{Eqn_nujs}
          \lim_{t \to \infty}  \frac{\Omega^{\mu_\epsilon}_{j}(t) }{t}
          \ \to \   \nu_j  := 1 - \rho_\epsilon( 1 - {P_B}_j)
          \text{ almost surely.}
        \end{align}}
      In other words, {\it the long run time average service rate seen by the
      $\t$-queue equals $\nu_j,$ which depends only on the blocking
      probability ${P_B}_j$ of the eager class, and not on the specific
      $\i$-sub-policy that produced that blocking probability.} Further under
      the SFJ limit, the service process seen by the $\t$-class becomes
      uniform. Specifically it follows from \cite[Theorem~1]{ANOR} that, as
      $\mue \to \infty$ in the SFJ limit,
\begin{align}
\nonumber
& \sup_{t \le W}  \left | \Omega^{\mu_\epsilon}_{j}(t)   - \nu_j t \right |  \to   0 \ a.s. \mbox{ for any finite $W,$ and } \\
& \Upsilon_j^{\mu_\epsilon}  \  \convas  \   \frac { B_\tau }{ \nu_j} \mbox{, both  for any  initial $\epsilon$-state, } \label{Eqn_Upsilon_conv}
\end{align}
  where $\Upsilon_j^{\mu_\epsilon}$ denotes the time required to
  finish $B_\tau$ amount of work using the service process
  $\Omega^{\mu_\epsilon}_{j}(\cdot)$.
  %These results were instrumental in deriving the static
  %Pseudo-conservation law.  
  This uniformity of the service process under SFJ limit enables a
  closed form characterization of the performance of the $\t$-class
  (see \cite[Theorems~2,3]{ANOR}). A key feature of the above results
  is a \emph{pseudo-conservation law} that expresses the performance
  of the tolerant class purely in terms of the blocking probability of
  the eager class, independent of the underlying $\i$-policy that
  produced the blocking probability. {\it We show an analogous
  pseudo-conservation for dynamic scheduling policies in this~paper.
}

%
Finally, we state the following assumption, which ensures that the
$\t$-queue remains stable under each
 %of the $\g$ sub-policies
 eager sub-policy, when
they are applied in a $\t$-static manner.

\begin{enumerate}[\textbf{B}.4] 
\item   {There exists $\delta_\tau  > 0$
  such that $\rho_{j} := \frac{ \lambda_\tau } { \mu_\tau \nu_j } < 1-\delta_\tau
  \mbox{ for all } j.$}
\end{enumerate}
Assumption~\textbf{B}.4 guarantees that the $\t$-system is stable in
the dynamic setting as well, as is shown by
Lemma~\ref{lemma:TauStability}.

\subsection{Some example models}
\label{sec:example_models}

We begin with the description of one example system that satisfies our
assumptions, in which the system capacity is not completely
transferred to one class at any time, but rather a fraction of it is
used by each $\i$-customer, whilst the left is utilized by one
$\t$-customer.
 
\noindent {\bf Capacity Division (CD-$(p,K)$) policy:} Each
$\i$-customer uses {\small $(1/K)$} part of the service capacity.  If
there are {\small $0 \le \ell \le K$} number of $\i$-customers
receiving service, then $\i$-customers are served at a net service
rate of {\small $({\ell}/{K})$}, while the $\t$-customer in service
(if any) is served at rate {\small $((K-\ell)/K)$}. This continues up
to {\small $K$} $\i$-customers, and any further $\i$-arrival departs
without service. Note that whenever an existing $\i$-customer departs,
{\it the service rate of the $\t$-customer gets increased by {\small
    $1/K$}}. Of course, $1/K \geq c_{\min}$ to satisfy Assumption
\textbf{A}.3. Further there is a prior admission control on
$\epsilon$-arrivals, they are admitted with probability $p$
independent of all other events.  This is the description of the
$\i$-sub-policy (as in \cite{Value,ANOR}). Now the top-level policy
varies the probability of admission $p$ (and possibly the maximum $\i$
occupancy $K$) based on the occupancy of the $\t$-queue.
 
We refer to the above $\i$-sub-policy as Capacity Division or briefly
as the CD-$(p,K)$ policy.  Note that the
CD-$(p,K)$ policy captures a multi-server setting for the eager class.
While the $\i$-scheduler need not be work conserving, the $\t$-class
uses all the left over capacity. %One can have other $\i$-sub-policies
% either designed by the system and/or influenced by the impatient
% response of the $\i$-class.  We describe a few~here.
Another example of an $\i$-sub-policy is the following.
    
\noindent {\bf Limited Processor Sharing (LPS-$(p,K)$) policy:} This
sub-policy, denoted by LPS-$(p,K)$ (as in \cite{Value,ANOR}), admits
an incoming eager job into the system with probability $p,$ so long as
the number of eager jobs already in service is less than or equal to
$K.$ The entire service capacity of the server is shared equally
between the eager jobs in service (i.e., each eager job gets served at
rate $1/\ell$ when there are $\ell$ jobs in service). As before,
$1/K \geq c_{\min}.$ Note that under the LPS-$(p,K)$ policy, the
tolerant class receives service only when there are no eager jobs in
the system.


\ignore{
  \noindent {\bf Balking and Reneging:} \sout{As a part of a
    particular $\i$-sub-policy, the system may allocate a certain
    number of servers (recall that the system may be viewed as
    multi-server from the standpoint of the eager class) and might
    allocate a certain amount of waiting space for $\i$-class.  The
    $\i$-customers might respond to the resources allocated based on
    their patience levels: a) an $\i$-customer may not enter the
    system depending upon the $\i$-number already in system according
    to some probabilistic rule, as in {\it balking} models; or b) may
    leave the system after waiting for an exponentially distributed
    patience time of rate $\alpha^{\mue},$ as in {\it reneging}
    models.  In the case of reneging, we will require that the
    parameter $\alpha^{\mue}$ scales linearly with $\mue$, i.e.,
    $\alpha^\mue = \alpha \mue$ for some $\alpha \in (0, \infty)$} (as
  in \cite{ANOR}).
}

  \noindent {\bf Top-level policies:} The top-level policy can chose
  any one of these sub-policies for any $\t$-state. For example, when
  the $\t$-occupancy is greater than a certain threshold $L,$ one may
  allocate fewer individual servers to $\i$-customers (using, for
  example, the CD policy), while one may allocate the entire capacity
  to the $\i$-class and may serve them in LPS mode when the
  $\t$-occupancy is smaller.

  \revision{ \noindent {\bf Switching between sub-policies:} Note that
    under Assumption {\bf A}.1, the transition between sub-policies,
    triggered by an arrival/departure in the $\t$ queue, is
    simplified---any left over eager jobs at the time of the
    transition are simply `flushed'. This simplifies our analysis, and
    as argued before, would have a negligible impact on the
    performance of both classes in the SFJ scaling regime. In
    practice, a more natural implementation strategy would be to
    simply let the new sub-policy dictate the scheduling of the
    existing eager jobs in the system after the change of the $\t$
    occupancy. This implies a possible re-assignment of the service
    rates of eager customers at the time of the $\t$-transition.
    % ongoing $\i$-busy period to complete under the same (previously
    % operational) sub-policy when there is an arrival/departure in
    % the
    % tolerant queue, and only apply the next sub-policy once the
    % eager
    % queue becomes empty.
    We prove formally that doing this does not affect our main results
    under the SFJ limit, for the special case of exponentially
    distributed eager job sizes (see Theorem~\ref{Thm_withou_A1}).}

  
  % \sout{Alternatively, one may allocate
  % (say) only one server to the $\i$-class when the $\t$-occupancy is
  % high (or low), which might lead to increased levels of
  % balking/reneging by the $\i$-class.}


  % In the following section, we characterize the performance of the
  % tolerant and eager classes under the SFJ limit.
